% GENERATED FILE. Edit canonical lessons, then run npm run generate:books.
% canonical-source-hash: fnv1a64:629aff48ba5fd178
% canonical-lessons: FR-C102-police, FR-C102-theatre, FR-C102-bibliotheque, FR-C102-boulangerie, FR-C102-mairie

\chapter{Police, Theatre, Library, Bakery, Town hall}
\label{ch:fr-police-theatre-library-bakery-town-hall}
\hlchaptermodality{102}

\begin{chapteropening}
\textbf{By the end of this chapter:} \emph{I can name the police, a theatre, a library, a bakery, the town hall.}

Complete the last lesson of chapter 102: I can name the police, a theatre, a library, a bakery, the town hall.
\end{chapteropening}

\section[la police]{la police — the police}

\label{lesson:FR-C102-police}



\begin{quote}
Before the new one: say the French for a village, then the French for the post office.
\end{quote}

\subsection*{You'll want to know: la police}
\textbf{la police} — \textquotedblleft{}the police\textquotedblright{}.

\begin{cousinweb}
\textbf{police} $\leftarrow$ Latin \emph{politia}, \textquotedblleft{}government, citizenship\textquotedblright{}, from Greek \emph{politeia}, built on \emph{polis}, \textquotedblleft{}a city\textquotedblright{} — a learned borrowing. (EtymDB 2.1)
\end{cousinweb}

\begin{grammarlens}[title={What you've built}]
One more word for this chapter.
\end{grammarlens}

\subsection*{Your turn}
\begin{itemize}
\raggedright
  \item \emph{Say it:} \emph{la police}
  \item \emph{Say it:} \emph{la police}, once more
  \item \emph{Say it:} say \emph{la police}
  \item \emph{Recall:} say the French for a village, then the French for the post office, then say \emph{la police} again
\end{itemize}

\begin{tcolorbox}[breakable,colback=teal!4,colframe=teal!35!black,title={Before you move on}]
What does la police mean? (\textquotedblleft{}The police\textquotedblright{}.) Say it once more.
\end{tcolorbox}
\section[le théâtre]{le théâtre — a theatre}

\label{lesson:FR-C102-theatre}



\begin{quote}
Before the new one: say the French for the post office, then the French for the police.
\end{quote}

\subsection*{You'll want to know: le théâtre}
\textbf{le théâtre} — \textquotedblleft{}a theatre\textquotedblright{}.

\begin{cousinweb}
\textbf{théâtre} $\leftarrow$ Latin \emph{theatrum}, from Greek \emph{theatron}, \textquotedblleft{}a place for viewing\textquotedblright{} — a learned borrowing. (EtymDB 2.1)
\end{cousinweb}

\begin{grammarlens}[title={What you've built}]
One more word for this chapter.
\end{grammarlens}

\subsection*{Your turn}
\begin{itemize}
\raggedright
  \item \emph{Say it:} \emph{le théâtre}
  \item \emph{Say it:} \emph{le théâtre}, once more
  \item \emph{Say it:} say \emph{le théâtre}
  \item \emph{Recall:} say the French for the post office, then the French for the police, then say \emph{le théâtre} again
\end{itemize}

\begin{tcolorbox}[breakable,colback=teal!4,colframe=teal!35!black,title={Before you move on}]
What does le théâtre mean? (\textquotedblleft{}A theatre\textquotedblright{}.) Say it once more.
\end{tcolorbox}
\section[la bibliothèque]{la bibliothèque — a library}

\label{lesson:FR-C102-bibliotheque}



\begin{quote}
Before the new one: say the French for the police, then the French for a theatre.
\end{quote}

\subsection*{You'll want to know: la bibliothèque}
\textbf{la bibliothèque} — \textquotedblleft{}a library\textquotedblright{}. A false friend: a bookshop is \textbf{la librairie}.

\begin{cousinweb}
\textbf{bibliothèque} $\leftarrow$ Latin \emph{bibliotheca}, from Greek \emph{bibliotheke}, \textquotedblleft{}a case for books\textquotedblright{}, built on \emph{biblion}, \textquotedblleft{}a book\textquotedblright{} — a learned borrowing. (EtymDB 2.1)
\end{cousinweb}

\begin{grammarlens}[title={What you've built}]
One more word for this chapter.
\end{grammarlens}

\subsection*{Your turn}
\begin{itemize}
\raggedright
  \item \emph{Say it:} \emph{la bibliothèque}
  \item \emph{Say it:} \emph{la bibliothèque}, once more
  \item \emph{Say it:} say \emph{la bibliothèque}
  \item \emph{Recall:} say the French for the police, then the French for a theatre, then say \emph{la bibliothèque} again
\end{itemize}

\begin{tcolorbox}[breakable,colback=teal!4,colframe=teal!35!black,title={Before you move on}]
What does la bibliothèque mean? (\textquotedblleft{}A library\textquotedblright{}.) Say it once more.
\end{tcolorbox}
\section[la boulangerie]{la boulangerie — a bakery}

\label{lesson:FR-C102-boulangerie}



\begin{quote}
Before the new one: say the French for a theatre, then the French for a library.
\end{quote}

\subsection*{You'll want to know: la boulangerie}
\textbf{la boulangerie} — \textquotedblleft{}a bakery\textquotedblright{}.

\begin{cousinweb}
\textbf{boulangerie} — built from the word for \textquotedblleft{}baker\textquotedblright{}, with \emph{-erie}, \textquotedblleft{}a place for\textquotedblright{}. The baker word came from Picard \emph{boulenc}, probably from Middle Dutch \emph{bolle}, \textquotedblleft{}a round loaf\textquotedblright{}.
\end{cousinweb}

\begin{grammarlens}[title={What you've built}]
One more word for this chapter.
\end{grammarlens}

\subsection*{Your turn}
\begin{itemize}
\raggedright
  \item \emph{Say it:} \emph{la boulangerie}
  \item \emph{Say it:} \emph{la boulangerie}, once more
  \item \emph{Say it:} say \emph{la boulangerie}
  \item \emph{Recall:} say the French for a theatre, then the French for a library, then say \emph{la boulangerie} again
\end{itemize}

\begin{tcolorbox}[breakable,colback=teal!4,colframe=teal!35!black,title={Before you move on}]
What does la boulangerie mean? (\textquotedblleft{}A bakery\textquotedblright{}.) Say it once more.
\end{tcolorbox}
\section[la mairie]{la mairie — the town hall}

\label{lesson:FR-C102-mairie}



\begin{quote}
Before the new one: say the French for a library, then the French for a bakery.
\end{quote}

\subsection*{You'll want to know: la mairie}
\textbf{la mairie} — \textquotedblleft{}the town hall\textquotedblright{}.

\begin{cousinweb}
\textbf{mairie} — built from the word for \textquotedblleft{}mayor\textquotedblright{}, with \emph{-ie}: the mayor's house. The mayor word $\leftarrow$ Latin \emph{maior}, \textquotedblleft{}greater\textquotedblright{} (Lewis \& Short).
\end{cousinweb}

\begin{grammarlens}[title={What you've built}]
One more word for this chapter.
\end{grammarlens}

\subsection*{Your turn}
\begin{itemize}
\raggedright
  \item \emph{Say it:} \emph{la mairie}
  \item \emph{Say it:} \emph{la mairie}, once more
  \item \emph{Say it:} say \emph{la mairie}
  \item \emph{Recall:} say the French for a library, then the French for a bakery, then say \emph{la mairie} again
\end{itemize}

\begin{tcolorbox}[breakable,colback=teal!4,colframe=teal!35!black,title={Before you move on}]
What does la mairie mean? (\textquotedblleft{}The town hall\textquotedblright{}.) Say it once more.
\end{tcolorbox}
